\section*{Chapter \ref{ch:iv2:macro-market-knowledge-and-pitches} --- Macro, Market Knowledge and Pitches}
\markboth{\protect\omTocMark{}{Answers: Macro and Pitches}}{\protect\omTocMark{}{Answers: Macro and Pitches}}

\begin{solution}{iq:iv2:macro-market-knowledge-and-pitches:1}
The bond falls about 7\%, and the stock about 20\%. With growth up half a point, the stock falls about 11\%.

The stock's price per dollar of dividend goes from $1/0.04 = 25$ to $1/0.05 = 20$. With $g = 4.5\%$, it's $1/0.045 = 22.2$.

So a long-duration stock can move more than a bond, and the reason for the rate move decides much of the outcome.

\iqwrong{Saying that rates up means stocks down, always. If rates rise on stronger growth, profits may more than offset it.}
\end{solution}

\begin{solution}{iq:iv2:macro-market-knowledge-and-pitches:2}
It's 2.5\%. That's the inflation rate at which the two bonds earn the same. It also holds an inflation risk premium and a liquidity gap.

$4.2 - 1.7 = 2.5$. If inflation averages more than that, the linked bond wins.

So it isn't purely a forecast. A risk premium would push it above expected inflation, and the linked bond's lower liquidity pushes it down.

\iqwrong{Reading it as the market's inflation forecast with no adjustment. It's close to one, but it can move for technical reasons in a stress.}
\end{solution}

\begin{solution}{iq:iv2:macro-market-knowledge-and-pitches:3}
Short rates rise, long rates usually rise less, the currency tends to rise, equities tend to fall, and gold is mixed. All this is relative to what was priced.

The chain: more inflation means tighter policy than expected. So front-end yields rise most, and the curve tends to flatten. Higher rates attract money, so the currency rises. Equities fall with the higher discount rate, and rate-sensitive sectors fall most. Gold has no coupon, so it falls with higher real yields, though it can gain from inflation fears.

To build your own answer, always say which channel you mean: expected policy, inflation compensation, or real rates. Then say what would break it: a central bank seen as behind the curve would lift long yields instead, and the curve would steepen.

\iqwrong{Giving a list of directions with no channel. Say why. A person can reason from a mechanism when the data is different, and a memorised list can't.}
\end{solution}

\begin{solution}{iq:iv2:macro-market-knowledge-and-pitches:4}
A 52\% chance of a cut was priced. A hold raises short rates by about 13 basis points. Equities tend to fall on that surprise. What moves prices is the gap from what was priced.

$(4.33 - 4.20)/0.25 = 0.52$. So after a hold, the expected rate rises from 4.20 to 4.33.

If 100\% of a hold had been priced, nothing would happen.

\iqwrong{Saying ``no change, so no reaction.'' No change in the rate can be a large change in expectations.}
\end{solution}

\begin{solution}{iq:iv2:macro-market-knowledge-and-pitches:5}
Reported revenue falls 8\%. Next I'd ask about costs: if more of them are in euros than sales are, profit falls by less. And I'd ask about hedging, and whether lower prices in dollars win them sales.

$0.8 \times 10\% = 8\%$. The other 20\% doesn't change.

\iqwrong{Assuming profit falls 8\% too. That holds only if costs sit in euros in the same share as sales. With more of the costs in euros, profit falls by less, and a US cost base hits much harder.}
\end{solution}

\begin{solution}{iq:iv2:macro-market-knowledge-and-pitches:6}
Learn a short list before each interview, with the date, and say so when you give it. If you don't know, say that, give the range you're sure of, and say how you'd look it up.

\textbf{How to prepare.} Make a one-page list and update it the morning of the interview. Rates: policy rate, two-year and ten-year government yields. Equities: main index and its recent move. Volatility: the main index's implied volatility. Currency: the main dollar pairs. Commodities: crude, natural gas, gold. Credit: an investment-grade and a high-yield spread. Note the last month's big move for each, and one reason for it.

\textbf{What to say.} ``As of this morning, about X. It's moved by about Y this month because of Z.'' If you're not sure: ``I'd say in the range of A to B, but I'd want to check.'' A wide range, honestly given, beats a confident wrong number.

\textbf{How to build your own.} Keep a source for each line, such as a central bank page or an exchange's data, so that you can say where it comes from.

\begin{dated}{2026-11}{Dates to know}\label{dat:iv2:macro-market-knowledge-and-pitches:cal}
In the US, the Fed schedules eight policy meetings a year, and the monthly jobs and consumer price reports come out at 8:30 a.m. Eastern time. Dates change, and a federal funding gap cancelled one jobs report in 2025. Look up the current calendar.
\end{dated}

\iqwrong{Bluffing a level. Interviewers can check it in a moment, and a wrong number delivered with confidence costs you more than saying that you don't know.}
\end{solution}

\begin{solution}{iq:iv2:macro-market-knowledge-and-pitches:7}
It's about \$238 million, or about 24\% of profit. With half hedged, it's about 12\%.

$\$10/42 = \$0.238$ a gallon, on 1 billion gallons. The hedged half is covered, so it's \$119 million.

\iqwrong{Forgetting fares. If competitors face the same cost, some of it reaches passengers. That's why the first-order number is a ceiling, but only if demand holds.}
\end{solution}

\begin{solution}{iq:iv2:macro-market-knowledge-and-pitches:8}
It implies about 3.875\% for the eight years. An inverted curve says the market expects lower short rates later, or a lower term premium. It doesn't say a recession is sure.

$((1.04)^{10}/(1.045)^2)^{1/8} - 1 = 3.875\%$. The quick version: $(10 \times 4.0 - 2 \times 4.5)/8 = 3.875\%$.

Past inversions have often come before recessions, in a small number of episodes, but the signal's timing is loose, and term premium changes can distort it.

\iqwrong{Treating correlation with past recessions as a causal law. It's a statement about priced policy paths, from a handful of cases.}
\end{solution}

\begin{solution}{iq:iv2:macro-market-knowledge-and-pitches:9}
The stock falls about 5\%: $+5\% - 0.5 \times 20 = -5\%$. They rise together when the higher yields come from stronger growth, and earnings rise more than the discount rate hurts.

The duration of 20 means 1 point of discount rate is 20\% of value. So 50 basis points is 10\%.

In a period with inflation well anchored, higher yields often mean better growth, so the correlation is positive. When inflation is the worry, higher yields mean tighter policy, and the correlation turns negative.

\iqwrong{Treating a fixed relation between the two. The sign depends on what's driving yields.}
\end{solution}

\begin{solution}{iq:iv2:macro-market-knowledge-and-pitches:10}
It gives 6\%, and 7.5\% with inflation a point higher.

$2 + 3 + 0.5 \times 1 + 0.5 \times 1 = 6$. With $\pi = 4$: $2 + 4 + 1 + 0.5 = 7.5$. So a 1-point rise in inflation raises the policy rate by 1.5 points, keeping real rates rising.

It can't tell you what's priced, and the output gap and $r^*$ can't be observed. A rule is a benchmark for how far policy is from past behaviour, and it doesn't decide anything.

\iqwrong{Saying the rule shows policy is wrong. Different estimates of $r^*$ and the gap can move the answer by a point or more.}
\end{solution}

\begin{solution}{iq:iv2:macro-market-knowledge-and-pitches:11}
It's about 9.2\% at $-0.3$, and about 10.8\% at $+0.3$.

The variance is $0.36 \times 0.0256 + 0.16 \times 0.0049 + 2 \times 0.24 \times \rho \times 0.0112 = 0.01 + 0.00538\rho$.

The sign depends on what drives the market. If growth shocks dominate, bonds rally when stocks fall, so the correlation is negative. If inflation and policy shocks dominate, higher rates hit both, so the correlation is positive.

\iqwrong{Assuming bonds are always a hedge. They were one over a long, low-inflation stretch, and that isn't a law.}
\end{solution}

\begin{solution}{iq:iv2:macro-market-knowledge-and-pitches:12}
It's about 5 basis points, a loss of \$50\,000. Yields might do nothing if it was already priced by other data, if revisions and wages are soft, or if the unemployment rate says otherwise.

$150\,000/30\,000 = 5$. And $5 \times \$10\,000$.

Payrolls have large revisions, and the household survey can disagree. So a beat can be undone within the release. The beta is an average across many regimes, and it changes when the Fed's focus changes.

\iqwrong{Trading the headline in isolation. Read the release: revisions, wages, unemployment and participation. And check what was priced after the earlier data.}
\end{solution}

\begin{solution}{iq:iv2:macro-market-knowledge-and-pitches:13}
The two-year falls about 15 basis points, and the ten-year about 5, so the curve steepens by 10, and the steepener makes about \$100\,000. The rest depends on why the bank was so bold.

The surprise is 25. So $0.6 \times 25 = 15$, and $0.2 \times 25 = 5$. The long two-year makes $15 \times \$10\,000$, and the short ten-year loses $5 \times \$10\,000$.

\textbf{The rest.} The currency tends to fall. Equities may rise on easier money. They may fall if the market reads the size as panic: ``what do they know that we don't?'' That's the point of an event question: name both readings, and say what would tell them apart.

\iqwrong{Using the full 50 as the shock. Only the surprise of 25 is news.}
\end{solution}

\begin{solution}{iq:iv2:macro-market-knowledge-and-pitches:14}
The dollar can fall about 3.8\% before you lose. The Sharpe is about 0.4. In an unwind, the carry doesn't protect you, and a sharp fall hits many holders at once.

$1.045 \times x/1.005 = 1$ gives $x = 0.9617$. That's $-3.83\%$. The carry of 4\% over a volatility of 10\% gives a Sharpe of 0.4.

A 1-standard-deviation fall of the dollar takes about 6\% off. Since many hold the same position with leverage, a fall forces others to sell, and that makes the return's tail fat: small gains for years, and a sudden large loss.

\iqwrong{Using the forward as a forecast, or saying the carry is free. A Sharpe of 0.4 is a modest return for a very negative skew.}
\end{solution}

\begin{solution}{iq:iv2:macro-market-knowledge-and-pitches:15}
It changes the supply of long bonds that private investors must hold, which affects their risk premium, and it signals the path of short rates. It's hard to measure, as both happen at once and the announcements are anticipated.

There are two channels. \textbf{Portfolio balance:} buying duration leaves less for private holders, so they accept less term premium. \textbf{Signalling:} a bank that's buying bonds is also unlikely to raise rates soon. Both lower long yields, and the opposite holds when it sells.

To build your answer: say what's a stock effect (how much it holds in total) and what's a flow effect (the pace). Say which you think matters. Say that estimates are usually from announcement windows, which catch only surprises and are themselves model-dependent.

\iqwrong{Saying ``they print money and so inflation follows.'' Reserves pile up at banks, and the link to broad money and to prices isn't mechanical.}
\end{solution}

\begin{solution}{iq:iv2:macro-market-knowledge-and-pitches:16}
The price falls about 5.5\%. The market cared about guidance, not the past quarter.

$1.05 \times 18/20 = 0.945$.

A beat is backward-looking, and the multiple is a price on the future. So the same print can give opposite reactions, depending on what's said about next year.

\iqwrong{Saying ``They beat, so it should go up.'' A price already held an expectation, for the numbers and for the guide.}
\end{solution}

\begin{solution}{iq:iv2:macro-market-knowledge-and-pitches:17}
Use this order: the one-line thesis, why it's mispriced, a number for what it's worth, the catalyst, the risks, and how you would size and exit. Be ready to defend each one.

\textbf{What must be in it.}
\begin{enumerate}
\item \textbf{Thesis in one line.} ``Long [invented company] since margins will recover faster than the market thinks.''
\item \textbf{The variant view.} What the market thinks and why it's wrong. No variant view means no pitch.
\item \textbf{A valuation number.} Your EPS against consensus, your multiple, your price. Bull, base and bear.
\item \textbf{A catalyst.} What makes others see it, and when.
\item \textbf{Risks.} Name the two that would hurt most, and say what you'd watch.
\item \textbf{Size and exit.} A stop or a thesis-break condition, and position size.
\end{enumerate}

\textbf{How to build your own.} Pick a company you can say something real about, in an industry you've looked at. Use the current numbers from its filings, and check them the morning of the interview. Have the multiple, the consensus and the main peers ready, with the date. Practise 2 minutes, then 5. Don't pitch a stock because it went up, or one everyone likes without saying what the market misses.

\iqwrong{Telling the company's story with no price. ``It's a great company'' isn't a pitch, as a great company at a high price can be a bad long. The interviewer wants what's in the price, and why you differ.}
\iqfollow{What if the interviewer says it's already priced in?}{Show what's priced, with a number: the multiple implies some growth rate. Then say why your number differs. If you can't, concede, and say what would change your mind.}
\iqfollow{And then a macro event hits?}{Run the event through your drivers: rates through the multiple, the currency through the sales mix, and input costs through the margin. Then say what you'd do to the position.}
\end{solution}

\begin{solution}{iq:iv2:macro-market-knowledge-and-pitches:18}
The expected return is about 9.4\%. Reward to risk is about 0.84 on the full cases, or about 2.5 to the stop. The position is about 8.3\% of the portfolio. A gap through the stop is what can go wrong.

\textbf{Cases.} Bull: $3.60 \times 18 = 64.80$, or $+29.6\%$. Base: $3.30 \times 16 = 52.80$, or $+5.6\%$. Bear: $2.70 \times 12 = 32.40$, or $-35.2\%$. The expected return is $0.5 \times 29.6 + 0.3 \times 5.6 + 0.2 \times (-35.2) = 9.44\%$.

\textbf{Reward to risk.} Bull gain over bear loss: $29.6/35.2 = 0.84$. Up to the stop, $29.6/12 = 2.47$.

\textbf{Size.} A \$6 stop on \$50 is 12\%. So $1\%/12\% = 8.3\%$ of the portfolio.

\textbf{The stop.} Earnings gaps go through stops. In the bear case, a gap down to 32.40 loses 35.2\%, or 2.9\% of the portfolio, about three times the plan. So if a report is before the catalyst, either cut the size or buy a put.

Also say how much of the 9.4\% beats what you'd make by owning the market with the same beta. That's the real edge.

\iqwrong{Sizing from the bull case. It's the loss that sets the size. Any stop that can be jumped over isn't a limit on what you lose.}
\end{solution}

\begin{solution}{iq:iv2:macro-market-knowledge-and-pitches:19}
The notionals are about \$52.6 million and \$12.5 million. Steepening of 25 basis points makes \$250\,000. Carry is about $-\$65\,000$ a year. VaR is about \$140\,000.

\textbf{Notionals.} $10\,000/(1.9 \times 0.0001) = \$52.6$ million, and $10\,000/(8 \times 0.0001) = \$12.5$ million.

\textbf{Profit.} The spread changes by 25, on \$10\,000: \$250\,000 for any parallel level. It's neutral to a parallel move.

\textbf{Carry.} Long two-year: $52.6 \times (4.0 - 4.1)\% = -\$52\,600$. Short ten-year: you pay 4.2\% and earn repo at 4.1\%, so $12.5 \times (-0.1)\% = -\$12\,500$. Together, about $-\$65\,000$ a year, so about \$5\,400 a month. Over three months, you need about 1.6 basis points of steepening to cover it. Rolling down the curve is left out here, and it should be added when the curve has a real slope at those points.

\textbf{Risk.} $6 \times \$10\,000 \times 2.326 = \$139\,600$.

\textbf{The pitch.} Why now: priced path against your view. Stop: say 15 basis points against you, a loss of \$150\,000. Don't forget the second notional is a quarter of the first, and it's the DV01 that matches.

\iqwrong{Using equal notionals. They'd make it a big bet on the ten-year, and not on the curve at all.}
\end{solution}

\begin{solution}{iq:iv2:macro-market-knowledge-and-pitches:20}
Nominal yields up 1 point: the nominal bond returns $-4\%$, and the linked one $+4.7\%$. Real yields up 1 point: the nominal bond returns $-4\%$, and the linked one $-3.3\%$. So it hedges inflation, not real rates.

\textbf{Nominal.} Its carry is 4, and the price move is $-8$. So $-4$ in both cases.

\textbf{Linked.} Its return is real yield 1.7, plus inflation 3, less duration on the real yield. So $4.7$ if the real yield is unchanged, and $4.7 - 8 = -3.3$ if it rises 1 point.

So when higher nominal yields come from break-evens, the linked bond is much better. When they come from real yields, it falls with the rest.

\iqwrong{Saying that an inflation-linked bond can't lose money in a high-inflation year. When the central bank responds and real yields rise, long-duration linked bonds can fall a lot.}
\end{solution}

\begin{solution}{iq:iv2:macro-market-knowledge-and-pitches:21}
Real yields moved, and that says tighter policy or better growth, with less expected inflation. Nominal bonds fall about 3.75\%, and linked bonds about 4.8\%. I'd rather own value stocks, as growth stocks have the longer duration.

$-7.5 \times 0.5 = -3.75$, and $-8 \times 0.6 = -4.8$.

The fall in break-evens points to a policy story: credibility is up, inflation fears are down, and real rates are higher. Not a pure growth story, in which break-evens would usually rise too.

Growth stocks have most of their value far off, so they're hit hardest by a real rate rise. Value has shorter cash flows, and it includes more cyclicals.

\iqwrong{Looking only at the nominal yield. A 50 basis point move can come from very different stories, with different winners. Always split it first.}
\end{solution}

\begin{solution}{iq:iv2:macro-market-knowledge-and-pitches:22}
A simple map: growth up and inflation down favours equities. Growth up and inflation up favours commodities and cyclicals. Growth down and inflation down favours bonds. Growth down and inflation up hurts almost everything except commodities and cash. The danger is that it hides what's priced.

\begin{center}
\begin{tabular}{lll}
\toprule
 & Inflation down & Inflation up \\
\midrule
Growth up & Equities win; bonds fine & Commodities win; equities mixed; bonds lose \\
Growth down & Bonds win; equities lose & Cash and commodities; both stocks and bonds lose \\
\bottomrule
\end{tabular}
\end{center}

\textbf{Why it's dangerous.} It sorts by level, while prices react to surprises. If high growth is priced, a merely good report disappoints. Assets also change: an energy-heavy index behaves differently from a tech-heavy one. And a few decades of data include few regimes.

\textbf{How to build your own.} Check the map against several past years, with the actual returns from a public source. Be ready to name one year that doesn't fit, and why.

\iqwrong{Reciting the table as though it were a law. Say ``tends to'', say that it's relative to what's priced, and say what evidence would change your mind.}
\end{solution}

\begin{solution}{iq:iv2:macro-market-knowledge-and-pitches:23}
The unhedged return is 4.5\%. Hedged, it's about 11.5\%. The hedge earns the rate difference as well as the euro's fall.

Unhedged: $1.10 \times 0.95 - 1 = 4.5\%$.

The forward is $F/S = 1.045/1.025 = 1.0195$. Selling euros forward on the starting value gains $1.0195 - 0.95 = 6.95\%$. So $4.5 + 6.95 = 11.45\%$.

Selling a low-rate currency forward means selling it at a premium. So the hedge earns the 2\% carry. If the euro doesn't move, you get 10\% plus about 2\%. The hedge covers only the starting value, so the 10\% profit remains exposed: the small residual here.

\iqwrong{Treating a hedge as a pure cost. It's a cost or a gain depending on the rate gap, and it reverses for a euro investor in the US.}
\end{solution}

\begin{solution}{iq:iv2:macro-market-knowledge-and-pitches:24}
They imply about 64\% for a cut at the first meeting and about 28\% for one at the second.

\textbf{First.} The month's average is half at 4.33 and half at the new rate $R$: $(4.33 + R)/2 = 4.25$. So $R = 4.17$, a priced fall of 16 basis points, or $16/25 = 64\%$.

\textbf{Second.} The month 2 future averages the rate after the second meeting: 4.10. That's 7 basis points below 4.17, or $7/25 = 28\%$.

So cumulatively, 23 basis points of cuts are priced by month 2.

\textbf{What to adjust.} Day counts: use the exact days before and after the meeting. Move sizes: a 50 might also be priced. Risk premium: futures aren't a pure forecast. The gap between the effective rate and the target's midpoint. And whether the contract settles on the average of the month, as this one does.

\iqwrong{Using 4.25 against 4.33 directly: $8/25 = 32\%$. That forgets that the meeting falls in the middle of the month, so half the average is already fixed.}
\end{solution}

\begin{solution}{iq:iv2:macro-market-knowledge-and-pitches:25}
The P\&L is $+\$1.4$ million, $+\$0.3$ million and $-\$2$ million. The mean is zero, and the standard deviation is about \$1.24 million. At 1\%, the straddle isn't worth buying on these numbers: its expected payoff is about 0.78\%. I'd cut the position by selling equity index futures or switching out of the two-year: the book is a single bet against high inflation.

\textbf{States.} P\&L in millions is $100r - 0.05\,\Delta y$:
\begin{itemize}
\item 0.2\%: $1.0 + 0.4 = +1.4$.
\item 0.3\%: $0.3 + 0 = +0.3$.
\item 0.4\%: $-1.5 - 0.5 = -2.0$.
\end{itemize}

\textbf{Mean.} $0.25 \times 1.4 + 0.5 \times 0.3 + 0.25 \times (-2.0) = 0$.

\textbf{Spread.} The variance is $0.25 \times 1.96 + 0.5 \times 0.09 + 0.25 \times 4 = 1.535$, and the root is 1.239.

\textbf{Straddle.} The expected absolute move is $0.25 \times 1 + 0.5 \times 0.3 + 0.25 \times 1.5 = 0.775\%$, against a cost of 1\%. So on your own numbers, you lose about 22 cents per dollar of premium. It would pay only if your map of moves is too small, or the distribution has more than these three states.

\textbf{Reading the position.} Stocks and two-years both lose in the bad state, so there's no diversification: they are the same bet. The two-year looks small next to \$100 million of stock, but it makes up 25\% of the worst case.

\textbf{Tonight.} The mean is zero, and the left tail is the bigger. So there's no reason to hold this size through the number. Decide your maximum loss, say \$1 million, and work back. Reducing the equities by 30\% takes the bad state to about $-\$1.55$ million. Closing the two-year as well takes it to $-\$1.05$ million. If the desk wants to keep the long, buy a put spread, and accept a cost that you can state.

\textbf{To the manager.} ``We're long one risk: inflation below consensus. A bad print costs about \$2 million, with a 25\% chance on my numbers. The rest of the distribution doesn't pay for that. I propose to cut by about a third, close the two-year, and re-enter after the release. And here's what would change my mind.'' Your map of states is a judgment: say so.

\iqwrong{Comparing each position's expected value and stopping. The two legs are correlated by the event, so a position that looks like a mix of equities and bonds is really a one-factor bet.}
\iqfollow{The print is 0.4\% and equities only fall 0.5\%. What do you think?}{The market moved less than your map. So either it was positioned for this, or it's reading something else in the release. Don't assume the remaining move is yours to fade. Check the two-year: if it agrees with your map, the equity reaction may lag, but say what you'd need to see.}
\end{solution}
